\paragraph{3.1. асимптотические обозначения}
\subparagraph{3.1.1}

Пусть $f(n)$ и $g(n)$ — асимптотически неотрицательные функции. Докажите с помощью базового определения $\Theta$-обозначений, что
\[
\max(f(n), g(n)) = \Theta(f(n) + g(n)).
\]

\textbf{Ответ:}

Заметим, что для любых неотрицательных функций $f(n)$ и $g(n)$ выполняются следующие неравенства:
\[
\max(f(n), g(n)) \leq f(n) + g(n),
\]
так как максимум не превышает сумму.

С другой стороны,
\[
f(n) \leq \max(f(n), g(n)), \quad g(n) \leq \max(f(n), g(n)),
\]
следовательно,
\[
f(n) + g(n) \leq 2 \cdot \max(f(n), g(n)).
\]

Отсюда получаем:
\[
\frac{1}{2}(f(n) + g(n)) \leq \max(f(n), g(n)) \leq f(n) + g(n).
\]

Таким образом, существуют положительные константы $c_1 = \frac{1}{2}$ и $c_2 = 1$, такие что
\[
c_1 (f(n) + g(n)) \leq \max(f(n), g(n)) \leq c_2 (f(n) + g(n)).
\]

По определению $\Theta$-обозначения получаем:
\[
\max(f(n), g(n)) = \Theta(f(n) + g(n)).
\]

\subparagraph{3.1.2}
Покажите, что для любых действительных констант $a$ и $b$, где $b > 0$, выполняется соотношение
\[
(n + a)^b = \Theta(n^b).
\]

\textbf{Ответ:}

Рассмотрим предел:
\[
\lim_{n \to \infty} \frac{(n + a)^b}{n^b} = \lim_{n \to \infty} \left(1 + \frac{a}{n}\right)^b = 1^b = 1.
\]

По определению предела, для $\varepsilon = \frac{1}{2}$ найдётся $N$, такое что для всех $n \ge N$:
\[
\frac{1}{2} \le \left(1 + \frac{a}{n}\right)^b \le \frac{3}{2}.
\]

Умножая на $n^b$, получаем:
\[
\frac{1}{2} n^b \le (n + a)^b \le \frac{3}{2} n^b.
\]

Следовательно, $(n + a)^b = \Theta(n^b)$ с константами $c_1 = \frac{1}{2}$, $c_2 = \frac{3}{2}$ (для всех $n \ge N$).
Для начальных $n$ константы можно при необходимости увеличить.


\subparagraph{3.1.3}
Поясните, почему утверждение «время работы алгоритма $A$ равно как минимум $O(n^2)$» лишено смысла.

\textbf{Ответ:}

$O(n^2)$ обозначает \textbf{верхнюю} границу: $T(n) \le c \cdot n^2$ для достаточно больших $n$.
Фраза «как минимум» требует \textbf{нижней} границы. Смешивать их некорректно, так как $O(n^2)$ включает функции,
растущие медленнее $n^2$ (например, $n$ или $\log n$), и утверждение «$T(n) \ge$ что-то из $O(n^2)$» не накладывает
содержательного ограничения снизу. Правильно: $\Omega(n^2)$ для нижней границы.

\subparagraph{3.1.4}
Справедливы ли соотношения
\[
2^{n+1} = O(2^n)
\]
и
\[
2^{2n} = O(2^n)?
\]

\textbf{Ответ:}

Первое справедливо:
\[
2^{n+1} = 2 \cdot 2^n = O(2^n),
\]
так как существует константа $c = 2$ и $n_0 \ge 1$, при которых $2 \cdot 2^n \le c \cdot 2^n$.

Второе не справедливо. Предположим противное: $2^{2n} = O(2^n)$.
Тогда по определению $\exists c, n_0: \forall n \ge n_0$ выполнено $2^{2n} \le c \cdot 2^n$.
Сокращая на $2^n$, получаем $2^n \le c$, что невозможно при $n \to \infty$.
Следовательно, соотношение не выполняется.

\subparagraph{3.1.5}
Докажите теорему 3.1.

\subparagraph{3.1.6}
Докажите, что время работы алгоритма равно $\Theta(g(n))$ тогда и только тогда, когда его время работы в наилучшем случае равно $O(g(n))$, а в наихудшем — $\Omega(g(n))$.

\subparagraph{3.1.7}
Докажите, что множество $o(g(n)) \cap \omega(g(n))$ является пустым.

\textbf{Ответ:}

Предположим противное: пусть существует функция $f(n) \in o(g(n)) \cap \omega(g(n))$.

По определению $f(n) = \omega(g(n))$: для любой константы $c > 0$ найдётся $n_1$, что при всех $n \ge n_1$ выполняется $f(n) > c \cdot g(n)$.

По определению $f(n) = o(g(n))$: для любой константы $c > 0$ найдётся $n_2$, что при всех $n \ge n_2$ выполняется $f(n) < c \cdot g(n)$.

Возьмём, например, $c = 1$. Тогда для достаточно больших $n$ (начиная с $\max(n_1, n_2)$) одновременно:
\[
f(n) > 1 \cdot g(n) \quad \text{и} \quad f(n) < 1 \cdot g(n).
\]
Получили противоречие: $f(n)$ не может быть одновременно строго больше и строго меньше $g(n)$.

Следовательно, такой функции не существует, и пересечение пусто.

\subparagraph{3.1.8}
Можно обобщить наши обозначения на случай двух параметров $n$ и $m$, которые могут возрастать до бесконечности по отдельности с разными скоростями. 
Для данной функции $g(n, m)$ обозначим как $O(g(n, m))$ множество функций
\[
O(g(n, m)) = \{ f(n, m) : \text{существуют положительные константы } c, n_0 \text{ и } m_0, \text{ такие, что } 0 \leq f(n, m) \leq c g(n, m) \}
\]
для всех $n \geq n_0$ или $m \geq m_0$.